% !TeX root = ../thesis.tex
\chapter{绪论}

\section{研究背景与问题提出}
\subsection{装备制造业知识平台建设背景}
随着制造业数字化转型持续推进，企业信息系统已从以 ERP、MES、PLM 为代表的流程管理系统，逐步扩展到面向生产知识、设备状态、工艺经验和运维规则的智能化知识平台。装备制造业企业在长期生产过程中沉淀了大量工艺规程、设备手册、故障记录、标准条文、维修案例和专家经验，这些内容既是企业核心知识资产，也是后续开展知识复用、辅助决策与智能分析的重要基础。然而，这些知识往往分散在不同业务系统、文档资料和人员经验中，表现为来源异构、语义关联分散、结构组织不统一和复用效率不高等问题\cite{manufacturingadoption2023,dskgsurvey2021}。从领域知识图谱研究的整体视角看，异构数据整合、模式设计、增量维护和跨系统复用仍是普遍存在的关键问题\cite{dskgsurvey2021}。

制造业知识平台建设的重点，是在现有业务系统之上形成统一的知识组织、管理和调用能力，使分散在不同来源中的文档信息能够逐步沉淀为可检索、可维护、可复用的结构化知识。对于装备制造业而言，这类平台尤其强调对设备、工艺、质量、故障和标准等对象之间关系的统一表达，并通过知识图谱、知识检索和模型服务等技术支撑知识汇聚与业务协同。知识图谱作为此类平台的重要组成部分，承担着将多源异构数据转化为结构化知识对象的任务，其质量在很大程度上依赖于上游知识抽取能力是否稳定、可扩展、可维护。

在实际工程中，装备制造业文本与通用互联网文本存在明显差异。首先，制造业文本中存在大量专业术语、型号缩写、复合参数和上下位概念，给实体边界识别和术语标准化带来困难。其次，工艺规程、故障分析和标准规范往往跨句甚至跨段表达关系，隐含因果、约束和前后条件较多，导致关系抽取难度明显高于一般新闻或开放领域文本。再次，企业平台对模型系统的要求并不止于抽取精度，还包括部署稳定性、资源可控性、版本可追踪性和结果可解释性等工程属性。因此，在制造业知识平台场景下，知识抽取任务已经不再是单纯的算法问题，而是算法、制品、服务和平台治理共同作用的工程问题。

\subsection{知识抽取在企业知识图谱中的关键作用}
知识图谱的核心在于将非结构化或半结构化信息转化为结构化实体、关系与属性。面向装备制造业知识平台建设，知识抽取主要承担三类作用。第一，它为知识图谱建设提供原始结构化输入，是设备对象、工艺参数、故障因果和标准约束等知识进入图谱的前置环节。第二，它为工艺规程、故障报告和标准规范等文本提供结构化处理能力，例如识别关键设备、故障部件、工序参数关联和约束条件等核心信息。第三，它为后续知识校核、知识入库和结果复核提供统一输出基础，使抽取结果能够进入平台流程而非停留在离线实验阶段。

如果没有稳定的知识抽取能力，知识图谱建设将严重依赖人工录入和规则维护，其成本高、速度慢且难以应对制造业知识持续更新的特点。尤其是在设备维护和工艺管理场景中，知识更新具有频繁、碎片化和场景相关性强的特征，单纯依赖静态规则难以长期维持效果。因此，引入可训练的知识抽取模型，并将其组织为可维护的模型库，是提高知识入库效率和结构化处理能力的必要路径。

另一方面，企业侧真实需求并不希望模型系统成为“研究孤岛”。企业更关心的是：不同模型在什么场景下适合使用；模型如何被平台调用；模型升级是否会影响既有业务；异常结果能否追踪到版本和配置；资源受限时如何选择更合适的部署方案。这些问题决定了知识抽取研究不能停留在离线实验层面，而应进一步进入统一服务化、性能评测和工程治理层面。

\begin{figure}[htbp]
  \centering
  \begin{tikzpicture}[
    box/.style={draw, rectangle, text width=2.45cm, minimum height=0.95cm, align=center, inner sep=3pt, font=\wuhao},
    widebox/.style={box, text width=4.2cm},
    arr/.style={-Latex, semithick}
  ]
    \node[box] (a) at (0,0) {业务文本来源\\工艺、故障、标准};
    \node[box] (b) at (3.2,0) {知识抽取模型库\\实体关系抽取};
    \node[box] (c) at (6.4,0) {知识图谱与知识库\\结构化沉淀};
    \node[box] (d) at (9.6,0) {平台应用层\\检索展示、入库、复核};
    \node[widebox] (e) at (4.8,-1.85) {回归与治理\\质检、回流、版本迭代};
    \draw[arr] (a) -- (b);
    \draw[arr] (b) -- (c);
    \draw[arr] (c) -- (d);
    \draw[arr] (d.south) |- (e.east);
    \draw[arr] (e.west) -| (b.south);
  \end{tikzpicture}
  \caption{知识抽取模型库在装备制造业知识平台中的位置}
  \label{fig:intro_knowledge_flow}
\end{figure}

\subsection{平台接入需求与模型库价值}
从平台接入流程看，装备制造业知识平台并不是直接调用零散模型代码，而是要求上游文本先经过结构化抽取、标签映射、结果校核和制品管理，再进入知识入库和业务使用环节。对于本文关注的知识抽取模块而言，关键不在于覆盖平台全部能力，而在于能否将抽取结果以统一接口、统一格式和统一版本管理方式接入既有流程。当抽取结果出现遗漏、关系歧义或配置不一致时，平台侧还需要能够追踪到具体模型版本和部署参数，并支持回归验证与替换。因此，知识抽取模型如果不能被工程化地组织和治理，就难以真正进入平台使用链路。

从平台建设实践看，企业知识平台项目早期往往先形成若干可演示的能力模块，再逐步整合为平台化能力。在这一阶段，知识抽取模块最常见的问题并不是“没有模型”，而是“模型彼此分散、运行方式不一、结果口径不同”。同一业务部门可能同时需要面向工艺文本的高精度抽取、面向故障报告的低时延识别，以及面向标准规范的批量结构化处理。如果模型之间缺乏统一的制品规范与部署接口，那么即使单个模型在实验上效果较好，也难以在平台中形成可复用、可比较、可维护的能力单元。围绕模型库开展系统建设，比围绕单一模型开展局部优化更符合企业知识平台项目的真实推进路径。

模型库的价值还体现在对知识更新机制的支撑上。装备制造业中的文档和经验知识具有长期累积、频繁修订和跨部门共享的特征，许多知识对象在不同时间节点会以不同文本形态出现。例如，同一设备既会出现在台账中，也会出现在工艺规程、故障报告和备件清单中；同一工艺参数既可能体现为标准阈值，也可能体现为异常案例中的偏离值。若没有多模型协同的抽取能力和统一的版本治理机制，这些跨文档、跨场景知识就很难被持续汇聚为可计算的企业知识资产。

因此，本文强调的模型库建设并不是简单罗列若干经典模型，而是尝试建立一套能够支持知识持续生产与持续治理的工程框架。只有当模型目录、制品结构、服务接口和性能指标口径被统一之后，平台才可能在不同业务时点、不同硬件条件和不同文本场景下稳定调用知识抽取能力。这种“把能力固化为平台资产”的思路，也是本文将模型工程化部署与性能优化纳入研究主线的重要原因。

\subsection{项目背景、需求分析与关键问题}
\par\textbf{项目来源与现实问题}
本文工作来源于装备制造业知识图谱知识库项目。该项目平台已经具备知识库与图谱应用的总体框架，但知识抽取部分在实际推进过程中存在明显短板。一方面，不同知识抽取模型的代码实现、训练脚本、输入输出格式和依赖环境彼此独立，导致模型难以统一管理，也难以快速比较和切换。另一方面，训练得到的模型往往停留在研究原型阶段，没有形成统一制品、统一服务和统一部署链路，因此难以真正接入平台体系。

从项目推进经验看，知识抽取模型进入企业平台大致需要经过四个环节：模型训练、模型制品化、服务化封装和平台部署验证。若其中任一环节缺失，模型就难以从实验结果转化为工程能力。此前项目中的主要问题恰恰在于第二、第三和第四环节衔接不足。具体表现为：模型版本缺少统一标识，难以追踪；模型导出形式不统一，难以直接部署；接口协议不统一，平台需要为不同模型分别适配；性能指标缺少统一口径，难以支撑工程选型。

\begin{figure}[htbp]
  \centering
  \begin{tikzpicture}[
    box/.style={draw, rectangle, text width=2.9cm, minimum height=0.95cm, align=center, inner sep=3pt, font=\wuhao},
    core/.style={draw, rectangle, text width=3.1cm, minimum height=1.0cm, align=center, inner sep=3pt, font=\wuhao},
    arr/.style={-Latex, semithick}
  ]
    \node[core] (c) at (0,0) {知识抽取模型\\平台化落地困难};
    \node[box] (a) at (-4.9,2.2) {模型实现分散\\训练环境不统一};
    \node[box] (b) at (0,2.7) {制品形式不统一\\难以标准部署};
    \node[box] (d) at (4.9,2.2) {推理接口不统一\\平台适配成本高};
    \node[box] (e) at (-3.5,-2.4) {性能口径不统一\\模型难比较};
    \node[box] (f) at (3.5,-2.4) {版本与监控缺失\\运维回归困难};
    \draw[arr] (a) -- (c);
    \draw[arr] (b) -- (c);
    \draw[arr] (d) -- (c);
    \draw[arr] (e) -- (c);
    \draw[arr] (f) -- (c);
  \end{tikzpicture}
  \caption{知识抽取模型从实验结果走向平台应用面临的主要问题}
  \label{fig:intro_problem_map}
\end{figure}

上述五个问题并非仅用于描述项目背景，而是构成后文章节展开的直接依据。针对“模型实现分散、训练环境不统一”问题，第三章通过多模型纳入、中文化适配和统一元数据组织降低模型之间的异构性；针对“制品形式不统一、推理接口不统一”问题，第三章和第四章进一步完成 ONNX 制品化、统一推理服务和部署链路设计；针对“性能口径不统一、模型难比较”问题，第二章与第四章建立统一评测指标体系和 Profile 驱动的 profiling 方法；针对“版本与监控缺失、运维回归困难”问题，第四章和第五章补充版本追踪、运行观测、基准测试与回滚验证机制。由此，绪论提出的问题在后续章节中均有对应的解决路径，论文整体结构也围绕这些问题逐项展开。

\par\textbf{业务需求分析}
结合制造业知识平台场景，本文认为知识抽取模型库需要满足如下需求：其一，在模型能力层面，需要覆盖关系抽取、联合抽取和复杂上下文建模等多类范式，以应对工艺文件、故障报告和标准规范等不同文本类型；其二，在工程层面，需要将不同模型统一组织为标准制品，并提供统一推理服务，降低平台接入成本；其三，在部署层面，需要通过显存、延迟、吞吐量和稳定性等指标支持模型比较与优化；其四，在系统层面，需要保留版本、日志、监控和回滚能力，使模型库能够被长期维护。

\begin{table}[htbp]
  \centering
  \caption{装备制造业知识抽取模型库的核心需求}
  \vspace{0.5em}\wuhao
  \begin{tabularx}{\textwidth}{l X X}
    \toprule
    需求类别 & 主要内容 & 对论文工作的要求 \\
    \midrule
    能力覆盖需求 & 支持工艺、故障、标准等文本场景下的实体与关系抽取 & 构建覆盖多类抽取范式的模型库 \\
    模型治理需求 & 不同模型应具备统一标识、统一制品和统一接口 & 完成模型制品化与服务化封装 \\
    部署评测需求 & 平台需比较不同模型在显存、延迟、吞吐量上的差异 & 建立统一部署评测指标体系 \\
    运维保障需求 & 模型升级、故障恢复和日志追踪应具备工程闭环 & 完成版本治理、监控与回滚设计 \\
    平台接入需求 & 外部系统能够按场景调用模型并获得统一结果结构 & 完成系统演示和应用验证 \\
    \bottomrule
  \end{tabularx}
  \label{tab:intro_requirement_expand}
\end{table}

\par\textbf{项目需求与验收要求}
若从利益相关方角度进一步细化需求，可以将模型库的直接使用者划分为算法研发人员、平台运维人员和业务使用人员三类。算法研发人员关注的是模型训练是否规范、标签与配置是否可复现、不同模型之间是否便于横向比较；平台运维人员更关心模型能否以统一制品形式上线、是否具备健康检查和回滚机制、资源占用是否可预测；业务使用人员则更加关注抽取结果是否符合业务语义、是否能快速返回、是否可以定位异常样例。三类角色关注重点不同，却共同构成了模型工程化建设的完整约束条件。

从工程验收视角看，装备制造业知识平台项目通常并不会以单一离线指标作为是否通过阶段验收的唯一依据。对于模型类成果，项目方往往需要同时确认四个问题：其一，模型是否已经形成标准化制品，能够被平台登记、上传和调用；其二，服务链路是否贯通，是否具备统一输入输出和异常处理能力；其三，部署后性能是否达到可接受范围，特别是显存占用、平均延迟、尾延迟和吞吐量是否满足典型场景要求；其四，系统是否具备日志追踪、版本回滚、样例验证和运行观测等工程支撑能力。企业验收更关注“是否可持续使用”，而非“某一次实验是否最好”。

进一步地，企业知识平台还要求模型库能够支持“逐步扩充”的演进模式。项目初期可能只接入少量模型和少量业务场景，但随着平台建设推进，新的文本类型、新的关系标签和新的部署节点会不断增加。如果论文提出的模型库方案只能在固定数据集和固定环境中运行，那么其工程价值就会显著受限。因此，绪论中的需求分析不仅要说明当前要解决什么问题，也要说明系统为何需要具备可扩展性、可比较性和可迁移性。这一点也是后续章节引入统一 ONNX 制品、统一推理接口与 Profile 模板的前提依据。

在这一背景下，本文将模型工程化部署理解为一种贯穿训练、导出、上线、监控和回归的全过程能力，而不是训练完成后的附加工作。只有当模型资产具备统一命名、统一配置、统一接口和统一监测口径时，平台侧才能真正把知识抽取模型看作稳定的服务组件。正因为如此，本文把需求分析延伸到工程验收层面，将“模型库建设”和“部署优化”同时视为研究对象，从而使论文结构更加贴近企业项目的真实推进逻辑。

进一步分析可知，企业对知识抽取模型的关注点与学术研究中的关注点并不完全一致。学术研究往往重点讨论公开数据集上的 Precision、Recall 和 F1，而企业更关心结果是否可被业务接受、接口是否稳定、资源是否可控以及版本是否可追踪。因此，本文在研究目标设置上，不单纯追求某一模型在某一数据集上的最佳成绩，而更重视“模型库是否构建完成”“服务链路是否打通”“部署是否具备优化能力”“系统是否具备演示与验证条件”这四类问题。

\par\textbf{关键问题定义}
围绕上述需求，本文将课题进一步归纳为三个关键问题。

第一，如何在多种经典知识抽取模型之间建立统一模型库。由于不同模型的输入结构、训练方式和输出语义存在差异，若直接将其堆叠在同一平台中，往往会导致训练流程复杂、接口适配成本高和版本治理困难。需要通过统一的元数据、统一的制品形式和统一的输出结构，降低不同模型之间的异构性。

第二，如何将训练结果转化为可部署、可比较、可回归的模型服务。模型训练完成后，若缺少统一的导出规范和服务封装方式，平台侧就难以稳定调用，也难以在不同模型之间进行横向比较。因此，需要将 ONNX 制品化、统一推理接口、Profile 配置和部署验证组织为连续链路。

第三，如何在企业场景中评价知识抽取模型的工程可用性。模型是否“可用”不仅取决于离线精度，还与显存占用、平均延迟、尾延迟、吞吐量、稳定性和回滚能力等因素有关。需要建立一套兼顾抽取质量与工程指标的评价口径，用于支撑后续平台选型与持续迭代。

\par
围绕上述三个关键问题，本文实际上试图回答一个更基础的工程命题，即“知识抽取能力如何从研究结果转化为平台资产”。这一命题要求论文不能仅对某一模型展开局部最优分析，而要对模型选型、中文适配、制品导出、服务封装、部署评测与系统验证等环节进行串联式论证。只有当这些环节被统一组织起来，知识抽取模型才能从实验代码演化为可注册、可发布、可治理的企业能力单元。
\par
本文的三个问题分别对应模型库构建、工程化部署和应用验证三个成果方向，其中模型库构建强调多模型统一组织，工程化部署强调服务稳定运行与性能优化，应用验证强调平台侧可调用、可分析和可验证。这样的组织方式使后续章节能够围绕同一主线展开，并形成方法、实现与验证相互衔接的结构。

\section{国内外研究现状与分析}
\subsection{知识抽取研究现状}
知识抽取研究经历了从流水线式任务拆分，到深度学习联合抽取，再到预训练语言模型驱动统一建模的发展过程\cite{resurvey2023,lmdsurvey2025}。以深度学习早期研究为例，实体识别与关系抽取常被拆分建模，分别采用序列标注、关系分类或联合标注方式实现\cite{lample2016,mahovy2016,chiu2016,miwabansal2016,zheng2017}。这类方法在任务定义清晰和结构约束显式方面具有一定优势，但对复杂上下文和长距离依赖的建模能力仍然有限。

Transformer 结构提出后，序列建模能力显著增强，为统一编码和联合抽取提供了更强基础\cite{vaswani2017}。BERT 及其变体进一步推动了知识抽取任务的性能提升，并带来了更强的表示与迁移能力\cite{bert2019,spanbert2020,tinybert2020}。在此基础上，不同研究围绕关系抽取和联合抽取形成了多种建模路线。例如，CasRel 将关系抽取拆解为主语识别与关系对象解码，结构相对清晰；SpERT 通过 span-based 范式增强实体边界识别能力；TPLinker 采用 token pair 标注处理重叠关系；PURE、PRGC、OneRel 和 CopyMTL 等方法则分别从任务解耦、阶段化建模、单步联合抽取和多任务复制机制等角度提升结构表达效果；DyGIE++ 与 GraphRel 则面向复杂依赖和图结构场景提供更强建模能力\cite{casrel2020,spert2019,tplinker2020,pure2021,prgc2021,onerel2022,copymtl2020,dygiepp2019,graphrel2019}。近年来，围绕关系特异性建模、语义增强以及任务分解策略，联合抽取方法仍在持续演进\cite{rsan2020,einet2022,decompose2024,mfie2024,singletable2024}。

总体来看，知识抽取方法已经形成较为丰富的模型谱系，但绝大多数研究仍以单模型效果为中心展开，对“多模型如何统一纳入同一工程体系”讨论较少。特别是在企业环境中，模型的选择标准不只由精度决定，还受资源约束、服务等级和运维要求影响。因此，构建模型库而非只比较单一模型，更符合企业知识平台场景的实际需要。

\subsection{制造业知识图谱与工程化研究现状}
在制造业领域，知识图谱研究逐渐从概念建模走向面向工艺、设备、产线和运维的综合应用。现有研究已在工艺规划知识组织、设备维护知识建模、故障诊断、数字孪生知识表示和制造过程分析等方向积累了较多成果\cite{processkgmethods2023,maintenancekg2023,faultkg2024,digitaltwinkg2024,prodlinekg2025,dtsystemkg2025}。近两年的研究还进一步关注制造知识图谱在故障诊断增强和智能诊断支持方面的演进趋势\cite{sensorkgicl2024,robotkg2025}。这些研究说明，制造业知识图谱在支撑经验沉淀、知识复用和智能决策方面具有现实价值。

然而，制造业知识图谱研究中对知识抽取模型的关注往往停留在“使用某类抽取方法完成局部实验”的层面，较少系统讨论面向平台落地的模型库组织方式。尤其是在企业级知识平台中，知识抽取并不是一次性的离线工作，而是需要长期运行、持续更新和可追踪治理的基础能力。因此，仅有图谱设计而缺乏稳定抽取模型体系，难以支撑知识库的持续演化。

\par\textbf{模型服务化与部署优化研究现状}
随着深度学习模型在工业界广泛应用，模型服务化与部署优化研究逐渐成为重要方向。Clipper 和 Clockwork 等代表性系统，从在线推理调度、资源配置和服务等级保障等角度讨论了模型服务系统的组织方式\cite{clipper2017,clockwork2020}。与此同时，面向跨框架模型交换、统一推理接口和性能 profiling 的工程实践不断成熟，围绕推理运行时、批处理、资源调度和性能观测的研究也为模型制品化与部署优化提供了较为明确的实现路径\cite{clipper2017,clockwork2020}。

这些成果为本文的工程实现提供了重要参考，但仍存在两个空缺。其一，现有推理基础设施研究主要面向通用模型服务，对知识抽取模型的结构特点关注不足。其二，现有部署优化研究较少与模型库建设结合，缺少“模型制品 + Profile 模板 + 统一服务 + 平台验证”的完整闭环。因此，本文尝试在制造业知识抽取模型库场景下，将这些工程能力组织为更具落地性的实现路径。

\par\textbf{面向工程落地的研究趋势}
近年来，知识抽取研究与模型服务化研究虽然分别快速发展，但二者之间正在出现更加明显的融合趋势。一方面，越来越多的工业场景研究开始关注模型输出是否保留证据片段、是否便于人工复核、是否能够支持增量更新等工程属性；另一方面，推理系统研究也不再只关注通用吞吐量提升，而开始强调针对不同任务形态设计更细粒度的配置模板和评测流程。对于知识抽取任务而言，这一趋势意味着研究重点正在从“是否能抽取”逐步扩展到“如何稳定抽取、如何高效抽取、如何长期维护抽取能力”。

这种融合趋势对装备制造业知识平台场景具有特殊意义。与开放问答、通用分类等任务相比，制造业知识抽取天然具有较强的证据追踪需求和版本追踪需求，下游往往需要知道某条关系来自哪一段文本、由哪个模型版本在何种 Profile 配置下产生。知识抽取模型的“生产过程”本身已经纳入平台治理范围。将这一治理过程纳入研究，可以更好体现工程型研究的价值，而不只是展示若干孤立的实验结果。

基于上述趋势，本文选择以“模型库构建 + 工程化部署 + 系统验证”作为整体切入路径。本研究既吸收知识抽取方法谱系中的代表模型，又吸收模型服务化和部署优化领域关于 profiling、指标治理和回归验证的思路，力图在装备制造业知识平台这一具体场景下形成较为完整的工程解决方案。这样既能够保持论文与已有研究的联系，也能够突出本文在落地性和系统性方面的差异化价值。

\subsection{现有研究不足与本文定位}
综合国内外研究现状，当前工作主要存在以下不足。

第一，知识抽取研究与企业平台部署之间存在明显断层。大量研究证明了某一模型在特定数据集上的有效性，但并未进一步回答模型如何被统一管理、统一服务和统一评测的问题。

第二，制造业知识图谱研究对抽取模型工程化能力关注不足。现有工作更多强调知识建模和业务价值分析，而较少围绕模型版本治理、运行监控、部署回滚和性能优化展开。

第三，面向多模型库的统一实现方法仍不充分。对于企业知识平台这类平台，真正需要的是“模型库能力”，即不同模型能够在统一接口和统一治理框架下按场景被选择和部署，而不是多个独立实验模型的简单并列。

\par\textbf{本研究的定位与切入点}
综合现有研究，当前真正缺少的并不是“还能否再提出一种知识抽取模型”，而是“如何把已有代表模型组织为可工程交付的模型库”。对于企业知识平台而言，模型的价值不在于单篇论文中的最优曲线，而在于它是否能够被稳定训练、统一导出、快速部署、按需选择并被系统持续维护。本文的研究定位正是在经典知识抽取模型基础上，建立一条适合装备制造业知识平台原型建设的工程化落地路径。

本文既强调研究对象的系统性和交付性，又保留了对模型方法谱系、中文适配策略和部署优化机制的系统论述。这种定位决定了全文既需要从文献层面对模型与服务化技术进行梳理，也需要从项目实践角度解释为何选择模型库、为何采用 ONNX 统一制品、为何使用 Profile 驱动部署优化。

本文的切入点并不试图覆盖企业知识平台中的所有模块，而是聚焦知识抽取模型库这一具有牵引作用的关键环节。其一，模型库直接决定了知识图谱上游结构化知识的质量与更新效率；其二，模型库天然连接算法训练、制品封装、部署上线和系统调用，是最能体现工程链路闭环的模块；其三，模型库具备较强的可扩展性，后续无论接入更多模型还是接入真实实验数据，都可以在本文原型基础上持续演化。因此，将其作为论文主体，既具有问题代表性，也具有后续扩展价值。

\par\textbf{研究问题的工程抽象}
从工程抽象角度看，本文所面对的问题可以归纳为“异构模型资产的标准化治理问题”。这里的“异构”既包括不同模型在训练框架、标签体系、输入长度和输出结构上的差异，也包括不同部署场景在硬件资源、时延目标和调用方式上的差异；而“标准化治理”则要求这些差异最终被收敛到统一的模型目录、统一的制品形式、统一的服务协议和统一的评价口径之中。只有完成这一层抽象，知识抽取模型库才可能脱离具体代码仓库和实验环境，进入平台化、资产化和可持续维护的状态。

工程抽象还意味着要明确成果的最终表现形态。对于本课题而言，真正的成果不只是若干训练权重，也不只是若干系统截图，而是由模型资产、部署模板、评测指标和演示系统共同组成的复合型成果集。其中，模型资产回答“能否抽取”，部署模板回答“如何运行”，评测指标回答“运行得怎样”，演示系统回答“能否被平台看见并被用户使用”。这种成果形态更接近工程项目中的可交付件集合。

此外，研究问题的工程抽象还有助于解释本文为什么不把“部署”和“优化”完全分割开来。对于模型工程化而言，部署并不是一个完成即结束的动作，而是伴随后续 profiling、配置更新、版本回归和故障回滚的持续过程。如果仅把部署理解为模型上线，把优化理解为独立实验，就会割裂工程实践中的真实链路。本文在绪论阶段明确这种工程抽象，目的是为第四章的章节组织提供理论依据，使后文的结构安排不是写作上的权宜之计，而是由问题本身的工程属性所决定。

最后，这种工程抽象也为论文后续版本迭代提供了稳定骨架。当前初稿 1.0 以合理推断生成数据和原型系统截图为主，重点验证结构、方法和链路；在后续初稿 2.0 中，完全可以将真实实验结果、真实性能曲线和更细致的案例分析替换到同一框架内，而无需改变论文主体逻辑。这说明本文绪论所建立的问题抽象不仅服务当前文本组织，也服务后续研究深化与成果完善。

\section{研究内容与技术路线}
\subsection{研究目标}
本文的总体目标不是提出一种新的知识抽取算法，而是围绕装备制造业知识平台场景，构建一套可训练、可封装、可部署、可监控、可比较、可选型的知识抽取模型库原型，并围绕部署性能开展系统化优化与工程验证，使其具备进入企业平台进一步集成的工程基础。该目标具体体现在以下三个方面。

一是模型层目标。建立覆盖多类经典知识抽取范式的模型库，并完成基于公开中文通用数据集的统一训练与中文场景适配，使模型具备面向制造业典型文本的迁移与验证基础。

二是工程层目标。将模型训练结果统一导出为标准制品，构建统一推理服务和本地化部署链路，并形成可复现的部署优化与性能对比方法。

三是系统层目标。基于 GPUStack 风格原型系统完成模型管理、部署管理、基准测试和结果验证演示，并通过制造业典型业务文本验证模型库的应用可行性。

\par\textbf{研究边界与数据来源说明}
本文对“面向装备制造业知识平台场景”的体现，主要体现在问题来源、系统需求、部署约束和应用验证四个层面，而不等同于全部训练过程均基于企业专有标注语料完成。受真实企业标注数据获取与构建条件限制，本文训练阶段采用公开中文通用数据集完成多模型统一训练、中文化适配和基础能力验证；在系统实现与应用验证阶段，则围绕工艺文件、故障报告和标准规范等装备制造业典型业务文本、知识库写入需求以及平台接入流程，验证模型库原型的工程适配能力。

\subsection{研究内容}
围绕上述目标，本文将研究内容组织为以下四项。

第一，完成知识抽取模型选型与模型库构建。本文纳入 CasRel、SpERT、TPLinker、PURE、PRGC、OneRel、CopyMTL、DyGIE++ 和 GraphRel 等代表模型，并在统一框架下完成模型资产整理。

第二，完成中文化适配和统一制品化。本文将可迁移模型统一替换为 bert-base-chinese 编码器，基于公开中文通用数据集重构分词、输入处理和标签映射流程，并将模型统一导出为 ONNX 制品，同时补充元数据、标签配置和默认部署参数。

第三，完成统一推理服务与部署优化方法设计。本文定义统一输入输出接口，引入 Profile 文件组织 batch、并发、长度、精度和后端等参数，并建立显存、吞吐量、平均延迟和稳定性等指标体系。

第四，完成系统实现与应用验证。本文以 GPUStack 风格前端为基础生成系统页面，围绕模型管理、部署、测试、监控和结果验证完成系统演示，并通过典型业务文本进行阶段性验证。

\begin{table}[htbp]
  \centering
  \caption{本文主要研究内容与预期产出}
  \vspace{0.5em}\wuhao
  \begin{tabularx}{\textwidth}{l X X}
    \toprule
    研究内容 & 核心工作 & 预期产出 \\
    \midrule
    模型库构建 & 模型选型、中文化适配、统一 ONNX 导出 & 面向制造业场景的知识抽取模型库原型 \\
    工程化部署 & 统一推理服务、Profile 设计、部署优化与版本治理 & 可部署、可比较、可回归的模型服务链路 \\
    系统实现 & GPUStack 风格原型界面、部署管理、基准测试、监控展示 & 可演示的系统原型与页面截图 \\
    应用验证 & 典型业务场景测试、结果分析、应用评价 & 面向平台集成的阶段性验证结论 \\
    \bottomrule
  \end{tabularx}
  \label{tab:intro_work_expand}
\end{table}

\begin{table}[htbp]
  \centering
  \caption{知识抽取模型库原型的能力映射}
  \vspace{0.5em}\wuhao
  \begin{tabularx}{\textwidth}{l X X}
    \toprule
    原型能力 & 论文中的具体体现 & 对应章节 \\
    \midrule
    可训练 & 统一纳入多种经典知识抽取模型，完成中文化适配、双语实验组织与统一训练流程迁移 & 第3章 \\
    可封装 & 将训练结果整理为 ONNX 模型文件、元数据、标签配置和默认 Profile 组成的标准制品包 & 第3章、第4章 \\
    可部署 & 设计统一推理接口与本地化部署链路，支持模型上传、实例创建、后端绑定和节点调度 & 第4章、第5章 \\
    可监控 & 保留日志、指标、任务历史、结果校验和导出能力，形成运行可观测与问题回溯闭环 & 第4章、第5章 \\
    可比较 & 通过统一评测指标、基准测试任务和 Profile 模板，对不同模型和不同部署配置进行横向比较 & 第4章、第5章 \\
    可选型 & 基于延迟、吞吐量、显存、冷启动时间和稳定性等指标，为不同业务场景提供模型与 Profile 推荐依据 & 第4章 \\
    \bottomrule
  \end{tabularx}
  \label{tab:intro_capability_mapping}
\end{table}

\par\textbf{“模型库原型”定位的具体体现}
表~\ref{tab:intro_capability_mapping} 说明，本文研究对象并不是某一个知识抽取模型的结构创新，而是以多模型资产为基础，围绕训练迁移、标准制品、统一服务、部署优化、运行监控和模型选型形成的一套工程化原型能力。其中，“可训练”解决模型能否进入统一实验与迁移流程的问题，“可封装”解决训练结果能否沉淀为稳定资产的问题，“可部署”解决模型能否真正以服务方式运行的问题，“可监控”解决上线后能否被持续观察与回溯的问题，“可比较”解决不同模型能否在统一口径下横向评估的问题，“可选型”则进一步回答在具体业务场景与资源约束下应选择何种模型与部署配置的问题。正是上述六类能力共同构成了本文所说的“知识抽取模型库原型”，也使论文能够从单一模型研究转向系统化工程验证。

\par\textbf{研究内容之间的闭环关系}
上述研究内容并非相互独立的四个任务，而是具有明显前后依赖关系的工程闭环。模型选型与模型库构建决定了后续可被部署和比较的对象集合；中文化适配与统一制品化决定了模型是否能够摆脱原始论文实现环境，进入稳定的工程运行状态；统一推理服务与部署优化决定了模型是否能够被平台高效调用，并形成面向业务场景的差异化部署策略；系统实现与应用验证则进一步检验前述工作是否真正转化为可展示、可操作、可分析的原型能力。四者前后衔接，共同构成本文的研究链路。

对于工程型论文而言，闭环结构的意义在于避免研究内容停留在“局部有效”。如果只完成模型训练而缺乏制品化，就无法证明模型具备接入平台的条件；如果只完成部署而没有统一性能口径，就无法说明模型为何适合某一业务场景；如果只有系统截图而没有模型能力支撑，则又会使论文偏离知识抽取模型库这一核心对象。因此，本文在绪论中明确给出研究内容之间的闭环关系，实质上是在为后续章节建立统一的评价框架。

此外，这种闭环组织还为后续 2.0 版本扩展真实实验数据预留了自然接口。当前阶段的推理生成数据主要承担验证方法与展示原型的作用，后续替换为真实实验结果时，无需改变论文总体结构，只需在既有模型、既有指标和既有系统流程基础上更新对应数据与分析结论即可。这说明本文的写作结构不仅服务当前初稿，也服务后续版本迭代，具备较好的持续完善空间。

\subsection{研究难点与技术路线}
\par\textbf{研究难点}
本文工作并非简单地将多个模型放入同一论文中，而是面临较强的工程复杂性。首先，不同知识抽取模型的输入表示、标签组织和训练逻辑差异较大，模型统一制品化存在较高门槛。其次，制造业文本场景复杂，长文本、复合术语和隐式关系使得模型迁移与部署评测不能直接套用公开数据集经验。再次，企业平台环境对显存占用、平均延迟、尾延迟和版本回滚等指标高度敏感，这使得模型部署必须与性能优化和运维治理同时考虑。

本文需要在“研究性”和“工程性”之间取得平衡。若过于偏重单模型实验，难以体现企业知识平台背景；若只讨论系统界面而忽略模型构建与优化方法，又难以完整呈现问题求解过程。因此，本文采用“模型库构建 + 工程化部署 + 系统实现验证”的综合组织方式，以确保研究主线完整。

\par\textbf{研究边界与基本假设}
本文聚焦于大模型流行之前的经典知识抽取模型及其工程化应用，不将通用大语言模型作为核心比较对象。这一边界与项目实际阶段相一致，也有助于将研究重点集中在模型库构建、部署治理和平台接入能力上。

在实验数据方面，本文当前阶段允许采用合理推断生成的阶段性数据，用于完成初稿结构与方法说明；后续 2.0 版本将以真实业务实验结果替换相应表格和结论。本文默认企业部署环境为受控 GPU 环境，平台支持本地模型文件上传和实例化部署，但不要求完整接入正式生产流程。上述假设保证了研究工作既贴近真实工程条件，又能够在论文阶段可控推进。

\par\textbf{技术路线}
基于前文问题分析与研究目标，本文技术路线可概括为：先从业务需求出发完成模型选型与模型库构建，再通过统一导出和统一推理服务打通制品链路，随后围绕 GPUStack 本地部署模式设计 Profile 驱动优化方法，最后通过系统演示与应用验证完成工程可行性说明。

\begin{figure}[htbp]
  \centering
  \begin{tikzpicture}[
    box/.style={draw, rectangle, text width=2.35cm, minimum height=0.95cm, align=center, inner sep=3pt, font=\wuhao},
    arr/.style={-Latex, semithick}
  ]
    \node[box] (a) at (0,0) {业务需求分析};
    \node[box] (b) at (3.0,0) {模型库构建};
    \node[box] (c) at (6.0,0) {统一服务与部署};
    \node[box] (d) at (9.0,0) {Profile 优化};
    \node[box] (e) at (12.0,0) {系统实现与验证};
    \draw[arr] (a) -- (b);
    \draw[arr] (b) -- (c);
    \draw[arr] (c) -- (d);
    \draw[arr] (d) -- (e);
  \end{tikzpicture}
  \caption{本文技术路线}
  \label{fig:intro_route_expand}
\end{figure}

\par\textbf{技术路线的阶段化特征}
从阶段划分看，本文技术路线可以进一步理解为“能力构建阶段、能力部署阶段、能力验证阶段”三个层次。第一阶段面向知识抽取模型本身，重点解决模型纳入、中文化适配和统一制品化问题，目标是形成可管理、可比较的模型资产；第二阶段面向模型服务运行，重点解决统一推理接口、GPUStack 本地部署、Profile 配置和性能优化问题，目标是形成可部署、可调优的服务能力；第三阶段面向系统演示与应用验证，重点解决模型管理展示、部署过程展示、基准测试展示和结果验证展示问题，目标是形成可证明、可汇报的工程原型。

这种阶段化路线体现了工程项目中常见的“先打通链路、再细化优化、最后固化展示”的推进规律。对于知识抽取模型库而言，如果在基础链路尚未贯通之前就直接追求极致性能，很容易导致优化目标失焦；反之，如果只完成最小可运行原型而缺乏后续优化和验证，则论文又难以体现研究深度。因此，本文将三类阶段明确区分，但又通过统一制品、统一接口和统一指标将其保持在同一主线之下，以保证技术路线既完整又具有可操作性。

另一方面，技术路线的阶段化特征也解释了为什么本文选择在第四章中将部署与优化放在同一章内而采用内部拆分的方式。部署解决“如何运行”，优化解决“如何更好运行”，二者共享同一套模型制品和同一套服务框架，将其强行拆成两个独立章节容易造成重复叙述。采用同章分节组织，则能够更自然地体现从基础部署到性能优化再到治理保障的连续逻辑。

\subsection{论文结构与主要贡献}
本文共分为五章。第一章为绪论，主要介绍研究背景、需求分析、国内外研究现状、研究目标、研究边界与技术路线；第二章介绍相关技术与工程基础；第三章说明知识抽取模型库构建方法；第四章围绕模型工程化部署与性能优化展开；第五章给出系统实现与应用验证。最后在结论部分总结全文工作并讨论局限与展望。

本文的主要贡献可归纳为以下四点。

第一，面向装备制造业知识平台场景构建了覆盖关系抽取与联合抽取任务的知识抽取模型库原型，将 CasRel、SpERT、TPLinker、PURE、PRGC、OneRel、CopyMTL、DyGIE++ 和 GraphRel 等代表模型统一纳入模型资产体系，使模型库具备可复用、可比较和可扩充的基础。

第二，完成了多模型的中文化适配与双语实验构建。本文将支持预训练语言模型编码器的模型统一替换为 bert-base-chinese，并同步修改分词器、输入处理、标签映射和训练流程，使经典英文实现能够基于公开中文通用数据集完成训练、收敛与基础能力验证。

第三，设计了统一模型封装与服务化部署流程。本文将训练后的模型统一导出为 ONNX 制品，补充元数据、标签配置和默认部署参数，构建统一推理服务，并接入轻量化 GPUStack 本地部署平台，形成平台外接式模型服务原型。

第四，面向部署阶段提出了配置驱动的性能 profiling 与工程化部署方案。本文通过 Profile 文件组织 batch size、并发数、最大长度、精度模式和推理后端等关键参数，并从显存占用、吞吐量、延迟、冷启动时间和稳定性等维度开展对比分析，为企业应用中的模型比较、模型选择与部署选型提供依据。
本文的研究定位既不是单纯追求某一知识抽取模型的最优精度，也不是仅做一个脱离模型能力的系统界面，而是尝试在装备制造业知识平台这一真实工程背景下，回答“经典知识抽取模型如何被组织为可工程交付的模型库”这一更具落地性的研究问题。本课题的难点并不只在于模型本身，而在于如何把模型能力稳定转化为企业知识平台可以接收、演示和持续维护的工程成果。对于装备制造业知识平台这类项目而言，真正具有长期价值的不是某一次抽取结果，而是能够不断吸收新数据、接纳新模型、支持新部署节点并保持结果口径一致的能力体系。后续章节将围绕“模型库能力形成”这一主线展开。
